\section{OWL2Gen: A Configurable OWL~2 Ontology Generator}
\label{sec:OWL2Gen_intro}

OWL2Bench provides a stable, domain-grounded setting in which OWL~2
reasoners can be evaluated while data scale and expressivity are varied.
However, some experiments require the ontology itself to become an
experimental variable. Reasoning behaviour can depend not only on data
size, but also on the OWL~2 constructs that occur, their frequencies,
the organization of the class hierarchy, and the way entities are
connected across axioms.

Real-world ontologies are valuable for evaluating reasoners under
representative modelling conditions, but they provide limited control over
these characteristics. Changing one structural or logical property while
keeping others fixed is generally difficult. A configurable synthetic
generator provides a complementary setting in which such characteristics
can be varied systematically.

OWL2Gen addresses this requirement by generating domain-independent OWL~2
ontologies from configurable parameters. Users can specify initial entity
quantities, OWL~2 constructs and their frequencies, hierarchy and nesting
parameters, entity-selection strategies, consistency verification,
serialization format, and random seed. The current implementation supports
OWL~2 DL and OWL~2 EL.

The purpose of OWL2Gen is therefore not to reproduce a particular
application domain, but to provide controlled ontology workloads in which
logical and structural characteristics can be varied for systematic
reasoner evaluation.


\subsection{Design Objectives}
\label{OWL2Gen_design_considerations}

OWL2Gen was developed according to the following objectives:

\begin{enumerate}

\item \textbf{Configurable generation.}
Users should be able to control the initial entity quantities, requested
OWL~2 constructs, construct frequencies, language setting, and selected
structural characteristics of an ontology.

\item \textbf{Generation from scratch.}
Ontologies should be generated without requiring an existing domain
ontology, allowing their structure and construct composition to be
controlled independently of a predefined model.

\item \textbf{OWL~2 compliance.}
Generated ontologies should conform to the selected OWL~2 DL or OWL~2 EL
setting. Constructs outside OWL~2 EL should not be selectable when that
setting is chosen.

\item \textbf{Consistency management.}
Generated ontologies should be checked for consistency, while an optional
mode should attempt to preserve consistency during generation.

\item \textbf{Structural variation.}
Different patterns of hierarchy formation, expression nesting, and entity
reuse should be supported so that ontologies with different structural
characteristics can be produced from comparable construct configurations.

\item \textbf{Characterization and inspection.}
Generated ontologies should be accompanied by structural metrics and an
interactive graph that allow users to inspect and compare the resulting
structures.

\item \textbf{Reproducibility.}
Users should be able to provide a random seed and retain the generation
configuration so that a generation run can be repeated.

\end{enumerate}


\vspace{1em}
\noindent
Table~\ref{tab:owl2gen_comparison} positions OWL2Gen with respect to the
configurable ontology and knowledge-graph generation approaches discussed
in Chapter~\ref{ch:litReview}. The comparison focuses on the capabilities
most relevant to the design objectives above.

\begin{table}[htbp]
\centering
\caption{Comparison of configurable ontology and knowledge-graph generation
approaches.}
\label{tab:owl2gen_comparison}
\small
\setlength{\tabcolsep}{3pt}
\resizebox{\linewidth}{!}{%
\begin{tabular}{lccccccc}
\toprule
\textbf{Approach} &
\textbf{\begin{tabular}[c]{@{}c@{}}OWL~2\\Support\end{tabular}} &
\textbf{\begin{tabular}[c]{@{}c@{}}From\\Scratch\end{tabular}} &
\textbf{\begin{tabular}[c]{@{}c@{}}Construct\\Selection\end{tabular}} &
\textbf{\begin{tabular}[c]{@{}c@{}}Construct\\Frequency\end{tabular}} &
\textbf{\begin{tabular}[c]{@{}c@{}}Structural\\Control\end{tabular}} &
\textbf{\begin{tabular}[c]{@{}c@{}}Consistency\\Checking\end{tabular}} &
\textbf{\begin{tabular}[c]{@{}c@{}}Structural\\Metrics\end{tabular}} \\
\midrule

OntoBench~\cite{Ontobench}
& \cmark
& \cmark
& \cmark
& \xmark
& \xmark
& \xmark
& \xmark \\

PyGraft~\cite{pygraft}
& Partial
& \cmark
& \xmark
& \xmark
& \cmark
& \cmark
& \xmark \\

DLCC~\cite{DLCC}
& \xmark
& \cmark
& \xmark
& \xmark
& \cmark
& \xmark
& \xmark \\

\textbf{OWL2Gen}
& \textbf{DL, EL}
& \cmark
& \cmark
& \cmark
& \cmark
& \cmark
& \cmark \\

\bottomrule
\end{tabular}%
}
\end{table}

The OWL2Gen framework presented here extends our previous work on
configurable ontology generation~\cite{singh2024owl2gen}. The extended
approach removes the dependency on a seed ontology and broadens the
generation process through additional OWL~2 construct support,
configurable entity-selection and nesting mechanisms, consistency
management, structural characterization, and visualization.


\subsection{Design and Implementation}


\subsubsection{User Interface and Workflow}
\label{sec:gui_OWL2Gen1}

OWL2Gen provides a graphical user interface that supports ontology
configuration, generation, verification, inspection, and download.
Figure~\ref{fig:owl2gen_configuration} shows the main stages of the
configuration workflow.

The landing page introduces OWL2Gen and summarizes the generation process
(Figure~\ref{fig:owl2gen_configuration}\subref{fig:owl2gen_about}).
The ontology-configuration page allows the user to select OWL~2 DL or
OWL~2 EL and optionally specify the initial numbers of classes, object
properties, data properties, and individuals
(Figure~\ref{fig:owl2gen_configuration}\subref{fig:owl2gen_constructs}).

The supported OWL~2 constructs are organized into collapsible categories
based on the OWL~2 Quick Reference Guide~\cite{owl2}. For each construct,
the user specifies the desired number of occurrences. Category-level
controls allow all constructs within a category to be selected or cleared
together. When OWL~2 EL is selected, constructs outside that profile are
disabled automatically.

Entity quantities and construct frequencies serve different purposes.
The entity fields determine the initial sizes of the entity pools used
during generation, whereas the construct fields specify the requested
frequencies of supported axioms and expressions. Keeping these controls
separate distinguishes vocabulary size from construct frequency.

The structural-configuration page provides control over hierarchy depth
and branching factor, maximum expression-nesting depth and per-level
nesting probability, and the entity-selection strategy
(Figure~\ref{fig:owl2gen_configuration}\subref{fig:owl2gen_structure}).
These parameters influence how the requested constructs are instantiated
and how entities are reused during generation.

The output-configuration page provides options for consistency
verification, verification timeout, ontology serialization, and random
seed (Figure~\ref{fig:owl2gen_configuration}\subref{fig:owl2gen_output}).
The random seed controls randomized choices made during generation,
allowing a run to be reproduced under the same configuration.

\begin{figure}[!ht]
    \centering

    \begin{subfigure}[t]{0.49\textwidth}
        \centering
        \includegraphics[width=\linewidth]{figures/OWL2Gen/OWL2Gen1.png}
        \caption{Landing page and workflow overview.}
        \label{fig:owl2gen_about}
    \end{subfigure}
    \hfill
    \begin{subfigure}[t]{0.49\textwidth}
        \centering
        \includegraphics[width=\linewidth]{figures/OWL2Gen/OWL2Gen2.png}
        \caption{Ontology setting, entities, and constructs.}
        \label{fig:owl2gen_constructs}
    \end{subfigure}

    \vspace{0.2em}

    \begin{subfigure}[t]{0.49\textwidth}
        \centering
        \includegraphics[width=\linewidth]{figures/OWL2Gen/OWL2Gen3.png}
        \caption{Structural configuration.}
        \label{fig:owl2gen_structure}
    \end{subfigure}
    \hfill
    \begin{subfigure}[t]{0.49\textwidth}
        \centering
        \includegraphics[width=\linewidth]{figures/OWL2Gen/OWL2Gen4.png}
        \caption{Generation and output configuration.}
        \label{fig:owl2gen_output}
    \end{subfigure}

    \caption{OWL2Gen user interface and configuration workflow.}
    \label{fig:owl2gen_configuration}
\end{figure}

After generation, the results page presents the consistency-verification
outcome and structural metrics. The generated ontology can also be
inspected through an interactive graph and downloaded in the selected
serialization format.

\begin{figure}[!ht]
\centering
\includegraphics[
    width=0.92\textwidth,
    height=0.62\textheight,
    keepaspectratio
]{figures/OWL2Gen/OWL2Gen5.png}
\caption{OWL2Gen results page showing the verification outcome, structural
metrics, and graph visualization.}
\label{fig:owl2gen_results}
\end{figure}


\subsubsection{Ontology Construction}
\label{sec:owl2gendl}

A generation request can be represented as the configuration tuple
\[
\mathcal{G}
=
\langle \mathbf{n},\mathbf{q},\mathbf{h},\mathbf{x},s,p,g\rangle .
\]

Here, $\mathbf{n}=(n_C,n_{OP},n_{DP},n_I)$ specifies the initial numbers
of classes, object properties, data properties, and individuals. These
quantities are optional. If omitted, OWL2Gen uses the larger of the
predefined default and the minimum required by the selected constructs.
The minimum is computed component-wise across the selected constructs,
since all constructs draw from shared entity pools. Explicitly specified
values below this minimum are rejected.

The construct-count mapping $\mathbf{q}$ assigns each selected construct
its requested number of occurrences. The hierarchy parameters
$\mathbf{h}=(d,b)$ specify target depth and branching factor, while the
nesting parameters $\mathbf{x}=(m,r)$ specify maximum nesting depth and
the probability of nesting at each level. The parameter $s$ identifies an
entity-selection strategy, $p\in\{\mathrm{DL},\mathrm{EL}\}$ denotes the
selected OWL~2 setting, and $g$ indicates whether
\emph{Guarantee consistency} is enabled.

Construction proceeds in five stages:

\begin{enumerate}[label=\roman*)]

\item \textbf{Configuration resolution and validation.}
Optional entity quantities are resolved and the submitted configuration
is validated against the selected constructs and OWL~2 setting.

\item \textbf{Entity-pool initialization.}
The class, object-property, data-property, and individual pools are
initialized according to $\mathbf{n}$. These pools may grow during
generation when additional entities are required.

\item \textbf{Axiom generation.}
The constructs requested by $\mathbf{q}$ are generated using the hierarchy,
nesting, and entity-selection parameters. Each candidate axiom is checked
by a local consistency guard before being added. When
\emph{Guarantee consistency} is enabled, consistency is additionally
checked after each construct batch.

\item \textbf{Ontology completion and verification.}
A completeness pass adds missing declaration axioms. Without
\emph{Guarantee consistency}, the completed ontology is checked once by
Openllet. With the option enabled, consistency checking is performed
incrementally during generation.

\item \textbf{Characterization and output.}
OWL2Gen computes the structural metrics and prepares the generated ontology
for inspection, visualization, and download.

\end{enumerate}


\paragraph{Entity representation and naming.}

All generated entities use the namespace
\texttt{http://owl2gendl.com/generated\#}. Each entity type has a fixed
prefix and a separate counter, producing names such as
\texttt{Class1}, \texttt{objectProperty1}, \texttt{dataProperty1}, and
\texttt{individual1}. These names identify entity type and order of
creation but do not encode domain meaning.

For example, generated axioms may include
\[
\texttt{Class10}\sqsubseteq
\exists\,\texttt{objectProperty2}.\texttt{Class3}
\]
and
\[
\texttt{Class106}(\texttt{individual1}).
\]

Nested expressions use entities from the same pools and may take forms
such as
\[
\texttt{Class105}\equiv
(\texttt{Class32}\sqcap\texttt{Class50})
\sqcup\texttt{Class73}.
\]
Here, an intersection is nested within a union. The recursive construction
of such expressions is controlled by $\mathbf{x}$.


\paragraph{Entity pools and selection strategies.}

Each entity type is managed by a separate pool. For every entity $e$, the
pool records a usage count $u(e)$ that is incremented whenever the entity
is selected. Let
\[
E_{\mathrm{available}}
=
\{e\mid u(e)<c\}
\]
denote the set of entities whose usage remains below the reuse limit $c$.
If this set is empty, a new entity is created. Otherwise, an entity is
selected according to $s$:

\begin{itemize}

\item \textbf{Uniform Random} assigns the same probability to every
available entity and introduces no preference based on prior use.

\item \textbf{Preferential} selects an entity with probability proportional
to $u(e)+1$. Entities that have already been used frequently are therefore
more likely to be selected again.

\item \textbf{Chain} selects the most recently used entity, creating a
sequential and concentrated pattern of reuse. The name describes the
selection procedure and does not imply that the complete ontology graph
forms a simple linear chain.

\item \textbf{Balanced Tree} selects an entity with the lowest usage count,
with ties resolved by insertion order. This distributes selections more
evenly across the pool.

\end{itemize}

If a construct requires distinct operands, entities already selected for
the current axiom are excluded. Selection is retried a bounded number of
times, after which a new entity is created if no suitable existing entity
is found. Consequently, the quantities in $\mathbf{n}$ specify initial
pool sizes rather than strict upper bounds on the completed ontology.


\paragraph{Accounting for requested construct counts.}

For every selected construct, OWL2Gen attempts to generate the number of
occurrences specified by $\mathbf{q}$. Nevertheless, the final number of
axioms need not equal
\[
\sum_k \mathbf{q}(k).
\]

The completeness pass may add declaration axioms for entities introduced
during construction, increasing the total. Conversely, when
\emph{Guarantee consistency} is enabled, a construct batch that cannot be
generated consistently within the retry limit may be omitted, decreasing
the total. Requested construct counts should therefore be interpreted as
generation targets, while the results page reports the quantities actually
produced.

\paragraph{Construct generation and ordering.}

OWL2Gen processes each construct requested in $\mathbf{q}$ and attempts to
generate its requested number of occurrences. The hierarchy parameters
$\mathbf{h}$ guide the depth and branching of the asserted class hierarchy,
while the nesting parameters $\mathbf{x}$ control the depth and probability
of nested class expressions. If these parameters are not specified, default
values are used. Nesting is determined locally for each generated expression.
These parameters therefore influence the generated structure without
determining its final measured values exactly.

Construct types are processed sequentially in a configurable order. By
default, this follows the construct catalog, which is organized according to
the OWL~2 Quick Reference Guide~\cite{owl2quickreference}. Because each
accepted axiom becomes part of the ontology used to generate and check
subsequent candidates, construct order can affect both later generation and
consistency checking.

OWL2Gen therefore allows the construct order to be varied rather than imposing
a single fixed strategy. Depending on the selected mode, constructs may follow
the default catalog order, be grouped by category, be randomized using a
request-specific seed, or follow an explicit user-supplied sequence. This
allows the effect of different orderings to be explored without assuming that
one ordering is preferable across all construct combinations and ontology
structures.


\paragraph{Local consistency guarding.}

Before a candidate axiom is added, OWL2Gen checks whether it conflicts with
constraints already established during generation. Common conflicts can be
identified from incompatible combinations of constructs over the same
entities. For example, an object property declared asymmetric cannot also
be declared symmetric, and classes declared disjoint cannot subsequently
be declared equivalent.

Some OWL~2 DL restrictions require additional checks involving previously
generated relationships:

\begin{itemize}

\item \textbf{Property simplicity.}
OWL~2 DL restricts the use of non-simple object properties in constructs
such as cardinality and self-restrictions. OWL2Gen tracks the relevant
property characteristics and prevents their use where a simple property
is required.

\item \textbf{Functionality and cardinality.}
Functional properties are prevented from being combined with restrictions
that require more than one value.

\item \textbf{Property-chain regularity.}
OWL2Gen tracks dependencies introduced by subproperty and property-chain
axioms and rejects additions that would violate OWL~2 DL regularity
restrictions.

\item \textbf{Identity and equivalence.}
Equivalence and identity relationships are considered transitively when
checking subsequent axioms. For example, if $a$ is declared identical to
$b$ and $b$ to $c$, a later assertion that $a$ and $c$ are different is
rejected.

\end{itemize}

When a candidate cannot satisfy the applicable constraints within the
allowed number of attempts, it is abandoned rather than adding an axiom
known to violate the tracked constraints.

These checks are intentionally local and do not constitute a complete
consistency procedure. Contradictions arising from interactions among
multiple axioms may only be detected through reasoning. The local guard
therefore prevents known conflicts during construction, while ontology
consistency is verified separately by a reasoner.


\paragraph{Reasoner-backed verification.}

OWL2Gen uses Openllet to verify generated ontologies. By default, the
completed ontology is checked once after construction. The reported outcome
is \textsc{consistent}, \textsc{inconsistent}, or
\textsc{not verified within budget}. The third outcome indicates that the
configured timeout expired and is not interpreted as evidence of
inconsistency.

When \emph{Guarantee consistency} is enabled, verification is performed
after the batch for each construct type is added. If a batch makes the
ontology inconsistent, it is removed and regenerated up to three times.
If all attempts fail, that construct batch is omitted and generation
continues.

This mode increases the likelihood of obtaining a consistent final
ontology, but may alter the requested construct distribution and increase
generation time. It should therefore be interpreted as a
consistency-aware generation mode rather than as an unconditional
guarantee that a reasoner will return a verdict.


\subsubsection{Characterization of Generated Ontologies}
\label{sec:metrics_owl2gendl}

The generation configuration specifies requested constructs and selected
structural targets, but it does not prescribe every entity choice made
during construction. Different random choices and entity-selection
strategies can therefore produce different relationships among the
generated entities. OWL2Gen reports structural metrics to make these
differences observable.

The metrics are descriptive rather than normative. Greater hierarchy
depth, average degree, clustering, or inferred-to-asserted ratio does not
necessarily indicate a better ontology. Their purpose is to complement
the construct configuration by describing how the generated axioms are
organized and connected.


\paragraph{Maximum nesting depth.}

Maximum nesting depth records the greatest recursive depth of a class or
data-range expression occurring within a single axiom. For example,
\[
A\sqsubseteq B
\]
is shallower than
\[
A\sqsubseteq\exists p.(B\sqcap\exists q.C).
\]

The metric describes syntactic expression depth. It does not by itself
measure logical complexity or predict reasoning time.


\paragraph{Hierarchy depth.}

Hierarchy depth is the length of the longest path in the asserted hierarchy
formed by \texttt{SubClassOf} relationships between named classes.
Anonymous class expressions are excluded. A larger value therefore
indicates a longer asserted sequence of class specialization.


\paragraph{Average branching factor.}

The average branching factor is the mean number of direct subclasses of
non-leaf classes in the asserted named-class hierarchy. It distinguishes
narrow hierarchies from broader structures in which superclass nodes have
several direct children.


\paragraph{Tangledness.}

Tangledness describes the prevalence of multiple inheritance in the
asserted named-class hierarchy. Let $C_H$ be the set of classes that
participate as subclasses in this hierarchy. It is calculated as
\[
T=
\frac{
|\{c\in C_H\mid |\operatorname{parents}(c)|>1\}|
}{
|C_H|
}.
\]

A value of zero indicates that no participating class has more than one
direct superclass, while larger values indicate that multiple inheritance
occurs more frequently.


\paragraph{Entity-connectivity graph.}

To measure how entities are shared across axioms, OWL2Gen constructs an
undirected, unweighted graph $G=(V,E)$. The vertex set $V$ contains the
named classes, object properties, data properties, and individuals
occurring in logical axioms. For each logical axiom, an edge is added
between every pair of distinct named entities occurring in that axiom.
Declaration and annotation axioms are excluded, repeated co-occurrences
produce a single edge, and self-loops are not included.

The average node degree is
\[
\overline{d}
=
\frac{1}{|V|}
\sum_{v\in V}\deg(v).
\]

It describes the average number of other entities with which an entity
co-occurs. A larger value indicates denser entity co-occurrence, but does
not necessarily imply greater semantic richness or reasoning difficulty.

The clustering coefficient measures the extent to which the neighbours of
an entity are also connected to one another. For a node $v$ with degree
$k_v\geq 2$,
\[
C_v
=
\frac{2t_v}{k_v(k_v-1)},
\]
where $t_v$ is the number of edges between neighbours of $v$. The reported
graph-level coefficient is the average of $C_v$ over eligible nodes.


\paragraph{Inferred-to-asserted subclass ratio.}

For an ontology whose consistency has been established, OWL2Gen measures
the extent to which classification expands the asserted named-class
hierarchy.

Let $S_{\mathrm{asserted}}$ be the set of explicitly asserted
\texttt{SubClassOf} pairs between named classes and
$S_{\mathrm{entailed}}$ the corresponding subclass pairs obtained after
classification. Trivial entailments, including reflexive subsumptions and
relationships involving \texttt{owl:Thing} or \texttt{owl:Nothing}, are
excluded. The ratio is

\[
R_{\mathrm{inf}}
=
\frac{
\left|
S_{\mathrm{entailed}}
\setminus
S_{\mathrm{asserted}}
\right|
}{
\left|
S_{\mathrm{asserted}}
\right|
}.
\]

For example, $R_{\mathrm{inf}}=2$ means that two additional named-class
subclass relationships are inferred for every explicitly asserted one.
The metric is not computed when consistency has not been established or
when $S_{\mathrm{asserted}}$ is empty. It characterizes expansion of the
named-class hierarchy only and does not include other forms of entailment.


\paragraph{Visualization of generated ontologies.}

In addition to the numerical metrics, OWL2Gen allows users to inspect a
generated ontology as an interactive graph. Classes, object properties,
data properties, and individuals are represented as typed nodes, while
relationships between named entities are represented as labelled edges.

This visualization is distinct from the entity-connectivity graph used for
average degree and clustering. The metric graph records only whether two
entities co-occur, whereas the visualization retains entity types,
construct labels, and direction where applicable.

To maintain responsiveness, the visualization displays at most 1,500
edges. If this limit is exceeded, the interface reports that the displayed
graph is partial; the complete ontology remains available through the
download function.


\subsection{Empirical Validation of OWL2Gen}
\label{sec:owl2gen_evaluation}

The evaluation demonstrates how OWL2Gen supports controlled variation in
ontology composition and structure. We examine variation across repeated
runs, the effect of entity-selection strategies, the scalability of
consistency-aware generation, and structural correspondence with real
ontologies. Rather than seeking a single best configuration, the goal
throughout is to show how much these outcomes vary as the configuration is
varied - across seeds, entity-selection strategies, and (with the reasoner
tier itself now independently configurable) the reasoner used for
consistency checking.

All experiments were conducted on a machine with an Intel Core i7-14650HX
processor (16 cores, 24 logical processors), 16~GB of memory, and Windows~11,
using Java~17 (Eclipse Temurin 17.0.19+10). No explicit JVM heap limit was
set; the default maximum heap was approximately 3.9~GB. Intermediate and
final checks used timeouts of 3 and 30 seconds, respectively, and each
generation run was limited to 120 seconds by the experimental script. These
bounds were chosen to keep the validation manageable while exercising
different OWL2Gen configurations. The four experiments below use a single,
fixed reasoner tier (Openllet~2.6.5) for both intermediate and final
consistency checks; a separate evaluation of reasoner behavior when
multiple tiers are configured is presented afterward.


\subsubsection{Structural Variation across Repeated Runs}

The first experiment examines whether repeated runs of the same configuration
produce structurally different ontologies. All 64 supported OWL~2 DL
constructs were requested five times each, except \texttt{SubClassOf}, which
was requested 40 times to provide sufficient scope for hierarchy formation,
giving 355 requested occurrences. Each request used 80 classes, 6 object
properties, 3 data properties, and 10 individuals, with a hierarchy target
depth of 5 and branching factor of 3. The default nesting parameters $m=2$
and $r=0.3$ were used. Uniform Random selection and \emph{Guarantee
consistency} were enabled, and seeds 1--10 were evaluated. All ten generated
ontologies were verified as consistent.

{\small
\begin{table}[htbp]
\centering
\caption{Structural variation across ten runs of the same configuration.}
\label{tab:metric-summary}
\small
\begin{tabular}{lrrrr}
\toprule
Metric & Mean & Std.\ dev. & Minimum & Maximum \\
\midrule
Final axiom count                 & 438.90 & 5.86  & 428   & 447   \\
Achieved hierarchy depth          & 4.30   & 0.64  & 3     & 5     \\
Achieved average branching factor & 1.21   & 0.07  & 1.11  & 1.33  \\
Tangledness                       & 0.185  & 0.048 & 0.11  & 0.25  \\
Achieved maximum nesting depth    & 2.00   & 0.00  & 2     & 2     \\
Average degree                    & 6.22   & 0.25  & 5.84  & 6.62  \\
Clustering coefficient            & 0.307  & 0.016 & 0.281 & 0.327 \\
Inferred-to-asserted ratio        & 21.85  & 4.02  & 17.17 & 29.82 \\
\bottomrule
\end{tabular}
\end{table}
}

Although the configuration is fixed, the resulting ontologies vary across
seeds. Average degree ranges from 5.84 to 6.62 and the inferred-to-asserted
ratio from 17.17 to 29.82, with variation also observed in hierarchy depth,
tangledness, clustering, and final axiom count.

The hierarchy results additionally show that $\mathbf{h}$ defines targets
rather than guaranteed outcomes. With 40 requested \texttt{SubClassOf}
axioms, achieved depth ranges from 3 to 5 and reaches the target of 5 in some
runs, while branching ranges from 1.11 to 1.33. The realized hierarchy
therefore depends on the structural targets, requested subclass axioms, and
entity-selection strategy. Maximum nesting depth is 2 in every run, matching
the configured $m=2$. Overall, a fixed OWL2Gen configuration can produce
different structural and inferred characteristics across seeds.


\subsubsection{Structural Effects of Entity-Selection Strategies}

The second experiment examines whether entity-selection strategy provides
additional control over ontology structure. All 64 supported OWL~2 DL
constructs were requested eight times, giving 512 requested occurrences.
Each strategy was evaluated over ten seeds with \emph{Guarantee consistency}
enabled.

\begin{table}[htbp]
\centering
\caption{Structural characteristics of the four entity-selection strategies
at 512 requested occurrences. Values are medians over runs verified as
consistent from ten seeds.}
\label{tab:strategy-comparison}
\small
\setlength{\tabcolsep}{2pt}
\begin{tabular}{lrrrrrrr}
\toprule
Strategy & Verified (/10) & Depth & Branch. & Tangled. &
Avg Deg. & Cluster. & $R_{\mathrm{inf}}$ \\
\midrule
Uniform Random & 10 & 0.50 & 0.50 & 0.00 & 7.16 & 0.326 & 13.53 \\
Chain          & 10 & 1.00 & 1.00 & 1.00 & 2.54 & 0.317 & 18.55 \\
Balanced Tree  & 3  & 1.00 & 1.00 & 0.00 & 7.38 & 0.306 & 11.81 \\
Preferential   & 10 & 1.50 & 1.00 & 0.07 & 6.21 & 0.514 & 10.70 \\
\bottomrule
\end{tabular}
\end{table}

The strategies produce distinct structural characteristics under the same
construct configuration. Chain has the lowest median average degree (2.54)
and highest tangledness, consistent with concentrated entity reuse. Balanced
Tree has the highest median average degree among its verified runs (7.38),
while Preferential produces the highest clustering coefficient (0.514).
The inferred-to-asserted ratio also varies, with the highest median for
Chain (18.55).

During these experiments, Balanced Tree exposed a reproducible
\texttt{NullPointerException} in Openllet~2.6.5, traced to a previously
documented defect in the \texttt{JoinCondition}
implementation.\footnote{\url{https://github.com/Galigator/openllet/issues/39}}
The documented correction was applied and the experiments repeated. Only
three of ten Balanced Tree runs then completed consistency verification; the
remaining runs terminated with other Openllet exceptions rather than an
inconsistent verdict. They are therefore treated as reasoner execution
failures, not as evidence of ontology inconsistency. Overall, entity-selection
strategy provides an additional mechanism for varying ontology structure
beyond construct frequencies alone.


\subsubsection{Scalability of Consistency-Aware Generation}

The third experiment examines \emph{Guarantee consistency} as the requested
ontology size increases. The 64 supported OWL~2 DL constructs were requested
two, five, and eight times, corresponding to 128, 320, and 512 occurrences.
Each configuration was evaluated over ten seeds for all four strategies.
Generation times are end-to-end and include reasoner-backed consistency
checks.

\begin{table}[htbp]
\centering
\caption{Runs verified as consistent and end-to-end generation time as the
requested ontology size increases. Times are medians with [minimum, maximum]
over verified runs.}
\label{tab:consistency-scaling}
\footnotesize
\setlength{\tabcolsep}{2.5pt}
\begin{tabular}{lrrrrrr}
\toprule
& \multicolumn{2}{c}{128} &
  \multicolumn{2}{c}{320} &
  \multicolumn{2}{c}{512} \\
\cmidrule(lr){2-3}
\cmidrule(lr){4-5}
\cmidrule(lr){6-7}
Strategy &
Verified & Time (s) &
Verified & Time (s) &
Verified & Time (s) \\
\midrule
Uniform Random
& 10/10 & 1.30 [1.08, 1.62]
& 10/10 & 2.27 [1.83, 4.09]
& 10/10 & 2.66 [1.92, 21.48] \\

Chain
& 10/10 & 1.24 [1.08, 31.55]
& 10/10 & 1.95 [1.84, 2.18]
& 10/10 & 21.28 [19.49, 23.82] \\

Balanced Tree
& 10/10 & 1.29 [1.18, 1.54]
& 6/10  & 29.93 [9.94, 71.80]
& 3/10  & 7.28 [2.98, 7.64] \\

Preferential
& 10/10 & 1.15 [0.31, 1.35]
& 10/10 & 1.68 [1.47, 1.90]
& 10/10 & 1.95 [1.75, 2.08] \\
\bottomrule
\end{tabular}
\end{table}

At 128 occurrences, all runs are verified and median generation times remain
low. At larger sizes, behavior becomes more strategy-dependent. Uniform
Random and Preferential remain verified across all seeds with relatively low
median times, whereas Chain reaches a median of 21.28\,s at 512 occurrences.

Balanced Tree shows the greatest variability. At 320 occurrences, six of ten
runs are verified, with times from 9.94\,s to 71.80\,s and a median of
29.93\,s. At 512, only three runs are verified, with times from 2.98\,s to
7.64\,s. The lower median should not be interpreted as improved scalability
because it is based on only three completed runs; the others terminate with
Openllet exceptions rather than inconsistent verdicts. Occasional long runs
also occur for Chain at 128 and Uniform Random at 512. Thus,
consistency-aware generation varies with workload size and entity-selection
strategy, while reasoner behavior can further affect verification.


\subsubsection{Structural Correspondence with Real Ontologies}
\label{sec:owl2gen_real_ontologies}

OWL2Gen is not intended to reconstruct existing ontologies. This experiment
instead examines whether configurations derived from real ontologies produce
structural characteristics that were not explicitly supplied to the
generator.

Six ontologies were considered: PROV-O~\cite{provo},
SSN~\cite{ssn-sosa}, SOSA~\cite{ssn-sosa,sosa-journal},
Pizza~\cite{pizza-ontology}, a GO slim ontology
~\cite{geneontology2000,geneontology}, and the W3C Wine
ontology~\cite{owlguide}. Observed entity quantities and supported construct
frequencies, together with hierarchy depth and average branching factor, were
used to configure OWL2Gen. In all six cases, the reference entity quantities
were sufficient for generation, so no additional entities were introduced
and the configured counts were retained. Tangledness, average degree,
clustering coefficient, and the inferred-to-asserted subclass ratio were held
out and measured only after generation.

For each reference-derived configuration, all four entity-selection
strategies were evaluated over ten seeds with \emph{Guarantee consistency}
enabled. Table~\ref{tab:real-ontology-comparison} compares reference values
with medians from verified runs. Bold values indicate held-out metrics for
which the reference value falls within the generated interquartile range.

{\scriptsize
\setlength{\tabcolsep}{3pt}
\renewcommand{\arraystretch}{1.1}
\setlength{\LTpre}{0pt}
\setlength{\LTpost}{0pt}

\begin{longtable}{llrrrrrrr}
\caption{Structural characteristics of the reference ontologies and the
corresponding OWL2Gen outputs. Generated values are medians over verified
runs from ten seeds.}
\label{tab:real-ontology-comparison}\\

\toprule
Reference & Strategy & Verified (/10) & Depth & Branch. & Tangled. &
Avg Deg. & Cluster. & $R_{\mathrm{inf}}$ \\
\midrule
\endfirsthead

\multicolumn{9}{c}{\tablename~\thetable\ continued} \\
\toprule
Reference & Strategy & Verified (/10) & Depth & Branch. & Tangled. &
Avg Deg. & Cluster. & $R_{\mathrm{inf}}$ \\
\midrule
\endhead

\midrule
\multicolumn{9}{r}{Continued on next page} \\
\endfoot

\bottomrule
\endlastfoot

PROV-O & \textit{Reference} & --- & 3 & 3.11 & 0.217 & 5.36 & 0.270 & 0.586 \\
       & Uniform Random & 10 & 4.50 & 1.35 & 0.35 & 5.06 & 0.15 & 5.94 \\
       & Chain & 0 & --- & --- & --- & --- & --- & --- \\
       & Balanced Tree & 10 & 3.50 & 1.26 & 0.27 & \textbf{5.20} & 0.12 & 4.30 \\
       & Preferential & 7 & 8.00 & 1.69 & 0.39 & \textbf{5.38} & 0.25 & 7.46 \\
\midrule

SSN & \textit{Reference} & --- & 1 & 1.75 & 0.167 & 5.04 & 0.510 & 0.000 \\
    & Uniform Random & 10 & 6.00 & 4.40 & 0.83 & 9.47 & 0.50 & 0.41 \\
    & Chain & 10 & 1.00 & 1.00 & 1.00 & 2.44 & 0.19 & 0.58 \\
    & Balanced Tree & 7 & 2.00 & 5.45 & 1.00 & 10.00 & 0.49 & 0.17 \\
    & Preferential & 10 & 8.00 & 3.82 & 0.72 & 9.79 & 0.56 & 0.13 \\
\midrule

SOSA & \textit{Reference} & --- & 0 & 0.00 & 0.000 & 1.03 & 0.000 & 0.000 \\
     & Uniform Random & 10 & 0.00 & 0.00 & \textbf{0.00} &
       1.35 & \textbf{0.00} & \textbf{0.00} \\
     & Chain & 10 & 0.00 & 0.00 & \textbf{0.00} &
       1.57 & \textbf{0.00} & \textbf{0.00} \\
     & Balanced Tree & 10 & 0.00 & 0.00 & \textbf{0.00} &
       1.08 & \textbf{0.00} & \textbf{0.00} \\
     & Preferential & 10 & 0.00 & 0.00 & \textbf{0.00} &
       1.56 & \textbf{0.00} & \textbf{0.00} \\
\midrule

Pizza & \textit{Reference} & --- & 6 & 4.00 & 0.012 & 16.30 & 0.622 & 2.336 \\
      & Uniform Random & 10 & 5.50 & 6.83 & 0.42 & 7.32 & 0.38 & 23.06 \\
      & Chain & 10 & 1.00 & 1.00 & 1.00 & 2.99 & 0.56 & 28.76 \\
      & Balanced Tree & 10 & 4.00 & 4.20 & 0.24 & 6.82 & 0.42 & 25.82 \\
      & Preferential & 10 & 0.00 & 0.00 & \textbf{0.00} &
        6.08 & 0.49 & 61.41 \\
\midrule

GO slim & \textit{Reference} & --- & 2 & 3.50 & 0.000 & 2.59 & 0.347 & 0.062 \\
        & Uniform Random & 10 & 4.00 & 1.25 & 0.21 & 3.37 & 0.28 & 0.57 \\
        & Chain & 10 & 1.00 & 1.00 & 1.00 & 2.31 & 0.18 & 0.87 \\
        & Balanced Tree & 10 & 1.00 & 1.00 & \textbf{0.00} &
          \textbf{2.59} & 0.21 & 0.02 \\
        & Preferential & 10 & 6.00 & 1.42 & 0.30 & 3.85 & 0.29 & 1.27 \\
\midrule

Wine & \textit{Reference} & --- & 3 & 1.50 & 0.071 & 16.59 & 0.485 & 3.683 \\
     & Uniform Random & 8 & 0.00 & 0.00 & 0.00 & 5.88 &
       \textbf{0.50} & \textbf{0.00} \\
     & Chain & 10 & 1.00 & 1.00 & 1.00 & 3.20 & 0.54 & 105.18 \\
     & Balanced Tree & 4 & 0.00 & 0.00 & 0.00 & 4.82 &
       \textbf{0.49} & 36.42 \\
     & Preferential & 10 & 0.00 & 0.00 & 0.00 & 6.29 &
       \textbf{0.54} & \textbf{23.18} \\

\end{longtable}
}

For PROV-O with Chain, all ten runs reached the final
consistency-verification timeout without a reasoner exception and are
therefore treated as unavailable rather than as structural mismatches.

Overall, correspondence on the held-out characteristics is limited. SOSA,
with its comparatively flat structure, shows the closest correspondence,
whereas Pizza and Wine show larger differences, particularly in average
degree and inferred subclass relationships. Even hierarchy depth and
branching, although supplied as targets, are not always realized exactly.
Matching entity quantities, construct frequencies, and hierarchy targets
therefore does not determine all structural properties of a reference
ontology; closer correspondence would require additional control over, for
example, entity reuse and co-occurrence across axioms.

\vspace{1em}
\noindent
Overall, the experiments show that OWL2Gen generates controlled but
structurally varied ontology workloads. Construct configuration,
entity-selection strategy, and randomization influence different structural
characteristics, while consistency-aware generation also depends on the
integrated reasoner. The experiments exposed reproducible Openllet failures,
illustrating the usefulness of configurable generation for reasoner testing
as well as benchmarking. The real-ontology comparison further shows that
controlled synthetic generation should not be interpreted as exact ontology
reconstruction.


\subsection{Limitations}
\label{sec:owl2gendl_limitations}

OWL2Gen generates domain-independent ontologies using generic entity names.
Its outputs therefore do not capture the semantics or modelling practices
of a particular application domain. The comparison with real ontologies
further shows that configurations derived from observed entity quantities,
construct frequencies, and hierarchy characteristics do not necessarily
reproduce other structural properties of the reference ontology. Supporting
closer structural or domain-specific correspondence would require
additional generation controls.

Requested entity and construct quantities are generation targets rather
than exact properties of the completed ontology. Entity quantities define
initial pool sizes and may grow during generation, while declaration axioms
may increase the final axiom count. Conversely, when
\emph{Guarantee consistency} is enabled, construct batches that cannot be
retained consistently may be omitted. Similarly, hierarchy depth and
branching factor are target characteristics and need not be achieved
exactly in every generated ontology.

Consistency checking currently relies on Openllet and a configurable
timeout. Consequently, \emph{Guarantee consistency} cannot ensure that the
reasoner will always return a verdict, and its practical scalability
depends partly on the integrated reasoner. Finally, the reported structural
metrics describe selected properties of generated ontologies and should
not be interpreted as measures of ontology quality or as general
predictors of reasoning difficulty.